% ═══════════════════════════════════════════════════════════════════════════════
%   FEG Prism — Kuratowski Calculus Engine Code Manual
%   XeLaTeX document · Styled after Modulo Synthesis Vol II
% ═══════════════════════════════════════════════════════════════════════════════
%
%   Compile with:
%     xelatex feg_prism.tex && xelatex feg_prism.tex
%
\documentclass[11pt, a4paper, openany]{book}

% ─── Geometry ─────────────────────────────────────────────────────────────────
\usepackage[
  top=2.5cm, bottom=2.5cm,
  left=3cm, right=3cm,
  headheight=14pt
]{geometry}

% ─── Fonts (XeLaTeX) ─────────────────────────────────────────────────────────
\usepackage{fontspec}
\setmainfont{LMRoman10}[
  BoldFont=LMRoman10-Bold,
  ItalicFont=LMRoman10-Italic,
  BoldItalicFont=LMRoman10-BoldItalic
]
\setmonofont{LMMonoLt10}[
  BoldFont=LMMonoLt10-Bold
]

% ─── Colors (Modulo Synthesis palette) ────────────────────────────────────────
\usepackage{xcolor}
\definecolor{structure}{HTML}{1A535C}   % Deep Teal
\definecolor{main}{HTML}{C6892A}        % Golden Ochre
\definecolor{second}{HTML}{2E4057}      % Slate Blue
\definecolor{third}{HTML}{8B2635}       % Burgundy
\definecolor{fourth}{HTML}{3A7D44}      % Forest Green
\definecolor{mainlight}{HTML}{FDF6E3}   % Solarized Base3
\definecolor{secondlight}{HTML}{EDF2F4}
\definecolor{thirdlight}{HTML}{F9EBEC}
\definecolor{fourthlight}{HTML}{EBF5EC}

% ─── Packages ─────────────────────────────────────────────────────────────────
\usepackage{amsmath,amssymb,amsthm,mathtools}
\usepackage{booktabs,enumitem,microtype,float,longtable,array}
\usepackage{fancyhdr,titlesec,titletoc}
\usepackage{tcolorbox}
\tcbuselibrary{skins,breakable,theorems}
\usepackage{etoolbox}
\usepackage{hyperref}
\hypersetup{
  colorlinks=true,
  linkcolor=structure, citecolor=second, urlcolor=second,
  bookmarksnumbered=true, pdfstartview=FitH,
  pdftitle={FEG Prism — Kuratowski Calculus Engine Code Manual},
  pdfauthor={Juan Pablo Silva Alvarado}
}

% ─── Macros ───────────────────────────────────────────────────────────────────
\newcommand{\lP}{\ell_{\mathrm{P}}}
\newcommand{\dS}{d_{\mathrm{S}}}
\newcommand{\MPl}{M_{\mathrm{Pl}}}
\renewcommand{\chaptername}{Chapter}
\newcommand{\code}[1]{\texttt{#1}}
\newcommand{\func}[1]{\texttt{\color{structure}#1}}
\newcommand{\modname}[1]{\texttt{\bfseries\color{main}#1}}
\newcommand{\flag}[1]{\texttt{-{}-#1}}
\newcommand{\file}[1]{\texttt{\itshape #1}}

% ─── Chapter/Section Styling ─────────────────────────────────────────────────
\titleformat{\chapter}[display]
  {\normalfont\Large\bfseries\color{structure}}
  {\small\color{structure!60}\MakeUppercase{\chaptertitlename}\ \thechapter}
  {4pt}{\Huge}
\titlespacing*{\chapter}{0pt}{-20pt}{20pt}

\titleformat{\section}
  {\normalfont\Large\bfseries\color{structure}}{\thesection}{0.6em}{}
\titlespacing*{\section}{0pt}{12pt}{6pt}

\titleformat{\subsection}
  {\normalfont\large\bfseries\color{structure!80!black}}{\thesubsection}{0.5em}{}
\titlespacing*{\subsection}{0pt}{8pt}{4pt}

% ─── Headers ──────────────────────────────────────────────────────────────────
\pagestyle{fancy}
\fancyhf{}
\fancyhead[LE,RO]{\small\thepage}
\fancyhead[LO]{\small\nouppercase{\rightmark}}
\fancyhead[RE]{\small\nouppercase{\leftmark}}
\renewcommand{\headrulewidth}{0.3pt}

% ─── Theorem Environments (from Vol II) ──────────────────────────────────────
\newcounter{thm}[chapter]
\renewcommand{\thethm}{\thechapter.\arabic{thm}}

% Definition Box
\newtcolorbox[use counter=thm]{definition}[1][]{%
  enhanced,breakable,
  colback=mainlight,colframe=main,
  colbacktitle=main,coltitle=white,
  fonttitle=\small\bfseries,
  title={Definition~\thethm\ifx&#1&\else: #1\fi},
  arc=1.5pt,boxrule=0.6pt,
  left=8pt,right=8pt,top=6pt,bottom=6pt,
  before skip=10pt,after skip=10pt
}

% Theorem Box
\newtcolorbox[use counter=thm]{theorem}[1][]{%
  enhanced,breakable,
  colback=secondlight,colframe=second,
  colbacktitle=second,coltitle=white,
  fonttitle=\small\bfseries,
  title={Theorem~\thethm\ifx&#1&\else: #1\fi},
  arc=1.5pt,boxrule=0.6pt,
  left=8pt,right=8pt,top=6pt,bottom=6pt,
  before skip=10pt,after skip=10pt
}

% Axiom Box
\newtcolorbox[use counter=thm]{axiom}[1][]{%
  enhanced,breakable,
  colback=thirdlight,colframe=third,
  colbacktitle=third,coltitle=white,
  fonttitle=\small\bfseries,
  title={Axiom~\thethm\ifx&#1&\else: #1\fi},
  arc=1.5pt,boxrule=0.6pt,
  left=8pt,right=8pt,top=6pt,bottom=6pt,
  before skip=10pt,after skip=10pt
}

% Result Box
\newtcolorbox{result}[1][]{%
  enhanced,breakable,
  colback=fourthlight,colframe=fourth,
  fontupper=\itshape,
  arc=1.5pt,boxrule=0.6pt,
  left=8pt,right=8pt,top=6pt,bottom=6pt,
  before skip=10pt,after skip=10pt,
  title={Key Result},
  coltitle=fourth!50!black, fonttitle=\small\bfseries,
  attach boxed title to top left={yshift=-2mm,xshift=2mm},
  boxed title style={colback=fourthlight,frame hidden},
  #1
}

% Introduction/Overview Box
\newenvironment{introduction}[1][Chapter Overview]{%
  \begin{tcolorbox}[
    enhanced,colback=structure!5,colframe=structure,
    colbacktitle=structure,coltitle=white,
    fonttitle=\small\bfseries,title={#1},
    arc=2pt,boxrule=0.7pt,
    left=8pt,right=8pt,top=6pt,bottom=6pt,
    before skip=12pt,after skip=12pt
  ]\itshape\begin{itemize}[leftmargin=*,itemsep=2pt,parsep=0pt]
}{%
  \end{itemize}\end{tcolorbox}
}

\setlist{nosep,leftmargin=*}
\setlist[itemize,1]{label={\color{structure}\textbullet}}
\setlist[enumerate,1]{label={\color{structure}\arabic*.}}

% ─── Title Page Macro (from Vol II) ──────────────────────────────────────────
\makeatletter
\newcommand{\subtitle}[1]{\gdef\@subtitle{#1}}
\newcommand{\extrainfo}[1]{\gdef\@extrainfo{#1}}
\gdef\@subtitle{}\gdef\@extrainfo{}
\renewcommand{\maketitle}{%
  \begin{titlepage}\centering
    \vspace*{4cm}
    {\Huge\bfseries\color{structure}\@title\par}\vspace{0.6cm}
    \ifx\@subtitle\empty\else{\Large\color{gray!80}\@subtitle\par}\fi
    \vspace{2cm}
    {\color{main}\rule{0.6\textwidth}{2pt}\par}\vspace{2cm}
    {\LARGE\@author\par}\vspace{1cm}
    {\large\@date\par}\vfill
    \ifx\@extrainfo\empty\else
      \begin{minipage}{0.8\textwidth}\centering\itshape\color{gray!60}\@extrainfo\end{minipage}\vspace{3cm}
    \fi
  \end{titlepage}
}
\makeatother

% ═══════════════════════════════════════════════════════════════════════════════
\title{FEG Prism}
\subtitle{Kuratowski Calculus Engine --- Code Manual}
\author{Juan Pablo Silva Alvarado}
\date{\today}
\extrainfo{%
  Part of Modulo Synthesis / Fractal Entropic Geometrodynamics\\[4pt]
  DOI: 10.5281/zenodo.18769707\\[4pt]
  ``The universe is strictly computable, finite, and structurally homeostatic.
  The rest is just counting.''}

\begin{document}
\maketitle
\frontmatter
\tableofcontents
\mainmatter

% ███████████████████████████████████████████████████████████████████████████████
\chapter{Theory}
% ███████████████████████████████████████████████████████████████████████████████

\begin{introduction}
  \item The physical framework: Fractal Entropic Geometrodynamics (FEG)
  \item Three Geometries of spacetime, matter, and logic
  \item Key invariants enforced by the simulation
\end{introduction}

\section{Fractal Entropic Geometrodynamics}

\textsc{FEG Prism} is a Monte Carlo engine implementing the Kuratowski Calculus
--- the mathematical framework underlying \emph{Fractal Entropic
Geometrodynamics} (FEG) by Juan Pablo Silva Alvarado
(DOI:~10.5281/zenodo.18769707).  The central claim of FEG is that four-dimensional
spacetime, the Standard Model particle spectrum, and quantum mechanics all emerge
from a single discrete structure: the Hasse diagram of a Poisson-sprinkled causal
set, constrained only by \textbf{planarity preservation} (Kuratowski's theorem).

The engine is named after its fundamental topological particle: the \textbf{Causal
Prism} $K_{2,N}$, a complete bipartite subgraph with two poles and $N \geq 3$
intermediaries.  These are the only non-trivial substructures permitted by the
triangle-free property of the transitive reduction.

\section{The Three Geometries}

\begin{definition}[Geometry of Order --- Phase 1]
  Poisson sprinkling into a 4D causal diamond produces $N$ events.  The
  \textbf{transitive reduction} (Hasse diagram) removes all redundant causal
  links: edge $(u,v)$ survives iff no intermediate $w$ satisfies $u \prec w
  \prec v$.  The result is a triangle-free directed acyclic graph (DAG) whose
  Lorentz invariance is emergent from the isotropy of the Poisson process.
\end{definition}

\begin{definition}[Geometry of Matter --- Phase 2]
  The \textbf{Causal Prism} $K_{2,N}$ is the unique non-trivial bipartite
  subgraph permitted in a triangle-free graph.  It consists of:
  \begin{itemize}
    \item Two \textbf{poles} (origin $u$, destination $v$) at graph-distance $\geq 2$
    \item $N \geq 3$ \textbf{intermediaries} $w_1, \ldots, w_N$ shared by both poles
    \item Intermediaries are mutually disconnected (guaranteed by triangle-freeness)
  \end{itemize}
  Topological mass is the integer $M = N$ (number of intermediaries).
  The mass hierarchy reduces to $N \in \{3, 4, 5+\}$ for electron, muon, tau.
\end{definition}

\begin{definition}[Geometry of Logic --- Computational Complexity]
  The 2-hop detection algorithm for Causal Prisms runs in $O(N)$ because the
  physics is local.  Hasse links have bounded proper time ($\tau \leq 8$);
  therefore $K_{2,N}$ poles must be within bounded causal distance.  The bounded
  degree (\code{MAX\_HASSE\_DEGREE} $= 15$) ensures $O(1)$ per-node work.
  This is not merely an optimisation --- it is a computational proof that
  \textbf{physical locality is an information principle}.
\end{definition}

\section{Key Invariants}

The following invariants are preserved throughout the simulation pipeline:

\begin{enumerate}
  \item \textbf{Triangle-freeness}: The Hasse diagram contains no 3-cliques.
        Any triangle $u \to w \to v$ with $u \to v$ would be transitively
        reduced.  This is a mathematical theorem, not an empirical check.
  \item \textbf{Planarity preservation}: $K_5$ threats are absorbed into the
        nearest pole via vertex contraction, enforcing Kuratowski's planarity
        criterion.  This is the discrete analogue of colour confinement.
  \item \textbf{Lorentz invariance}: No spatial cutoff is applied during Hasse
        construction.  The light cone search at time offset $\Delta t$ uses
        radius $\Delta t + 1$ (natural units $c = 1$), preserving exact
        causal structure.
  \item \textbf{Strict finitism}: All accumulations use \code{u64} integers.
        The single \code{f64} division occurs only at the final return.
        No floating-point drift across ensemble realisations.
  \item \textbf{Deterministic reproducibility}: Given the same seed, the engine
        produces identical output on the same platform.  Seeds are derived via
        \code{StdRng::seed\_from\_u64(seed\_base + k)} for realisation $k$.
\end{enumerate}


% ███████████████████████████████████████████████████████████████████████████████
\chapter{Architecture}
% ███████████████████████████████████████████████████████████████████████████████

\begin{introduction}
  \item Six-phase simulation pipeline
  \item Module organisation and data flow
  \item Execution modes: in-memory vs.\ streaming
\end{introduction}

\section{Pipeline Overview}

The simulation proceeds through six phases, each corresponding to a Rust module:

\begin{verbatim}
  ┌─────────────────────────────────────────────────────────┐
  │  Phase 1: Sprinkle + Hasse       (phase1/)              │
  │    Poisson sprinkling → Transitive Reduction (CSR)       │
  ├─────────────────────────────────────────────────────────┤
  │  Phase 2: Kuratowski Contraction  (phase2/)              │
  │    K_{2,N} detection → K_5 threat absorption → Defect    │
  ├─────────────────────────────────────────────────────────┤
  │  Phase 3: Spectral Observables    (phase3/)              │
  │    Random walkers → d_S(t) → Causal Flux                 │
  ├─────────────────────────────────────────────────────────┤
  │  Measurements M1–M14              (measure/)             │
  │    M1–M10: Mass, half-life, modulo, vacuum pol.,         │
  │    electroweak, decoherence, neutrino, PMNS, Higgs,      │
  │    Lagrangian; M11–M14: Hausdorff, zigzag, collider,     │
  │    mass formula                                           │
  ├─────────────────────────────────────────────────────────┤
  │  Phase 4: Output                  (output/)              │
  │    CSV serialisation + provenance stamps                  │
  ├─────────────────────────────────────────────────────────┤
  │  Ensemble Runner                  (ensemble/)            │
  │    Adaptive Welford convergence + checkpointing          │
  └─────────────────────────────────────────────────────────┘
\end{verbatim}

\section{Module Map}

\begin{longtable}{lp{9cm}}
  \toprule
  \textbf{Module} & \textbf{Responsibility} \\
  \midrule
  \modname{main.rs}        & CLI parsing, flag dispatch, provenance boot \\
  \modname{config.rs}      & TOML configuration, RAM-aware mode selection \\
  \modname{provenance.rs}  & Five-layer SHA-256 provenance system \\
  \modname{phase1/}        & Poisson sprinkling (\file{sprinkle.rs}) and Hasse
                             construction (\file{hasse.rs}) \\
  \modname{phase2/}        & Causal Prism detection (\file{defect.rs}),
                             topology summary (\file{topology.rs}) \\
  \modname{phase3/}        & Spectral dimension (\file{spectral.rs}),
                             random walkers (\file{walker.rs}),
                             causal flux (\file{flux.rs}) \\
  \modname{graph/}         & Type-safe CSR (\file{csr.rs}),
                             grid quantisation (\file{grid.rs}) \\
  \modname{ensemble/}      & Adaptive batching (\file{runner.rs}),
                             Welford averaging (\file{averaging.rs}),
                             checkpointing (\file{checkpoint.rs}) \\
  \modname{measure/}       & Fourteen observer measurements M1--M14, in
                             \file{m01\_traversal.rs} through \file{m14\_mass\_formula.rs} \\
  \modname{output/}        & CSV writer (\file{csv.rs}),
                             human-readable summary (\file{summary.rs}) \\
  \modname{anim\_export.rs} & Binary topology slice for Manim animation \\
  \bottomrule
\end{longtable}

\section{Execution Modes}

\begin{definition}[In-Memory Mode]
  All $N$ points and the full CSR adjacency are held in RAM.  This is the fast
  path.  Memory estimate: $N \times 2000$ bytes for $N \leq 50{,}000$;
  $N \times 500$ bytes for $N > 50{,}000$ (Hasse degree bounded by 15).
  Selected by \flag{inmemory} or recommended automatically when estimated memory
  fits within $70\%$ of available RAM.
\end{definition}

\begin{definition}[Streaming Mode]
  For extreme scales ($N > 10^7$), the simulation uses time-slab sprinkles and a
  sliding window.  Only the active window is held in RAM.  Estimate:
  $N \times 13 + 100\,\text{MB}$ overhead.  Selected by \flag{stream}.
\end{definition}

The mode is auto-detected by \func{config::recommend\_mode()}, which queries
\func{sysinfo::System} for available memory and applies the safety fraction
from \file{causal\_set.toml}.


% ███████████████████████████████████████████████████████████████████████████████
\chapter{Phase 1 --- Poisson Sprinkling}
% ███████████████████████████████████████████████████████████████████████████████

\begin{introduction}
  \item Sprinkling $N$ events into a 4D causal diamond
  \item Volume formula and density normalisation
  \item Streaming chunks for sliding-window mode
\end{introduction}

\section{The Causal Diamond}

Source: \file{phase1/sprinkle.rs}

The simulation domain is a 4D causal diamond: the intersection of the future
light cone of a past tip $p^-$ and the past light cone of a future tip $p^+$.
In Cartesian coordinates, the constraint is
\[
  |t| + r \leq T/2, \qquad r = \sqrt{x^2 + y^2 + z^2}.
\]
The diamond has 4-volume $V = \pi T^4 / 24$.  To normalise the Poisson density
to $\rho \approx 1$, we set
\[
  T = \left(\frac{24N}{\pi}\right)^{1/4}.
\]

\section{Rejection Sampling}

The function \func{sprinkle(n, rng)} draws uniform random points in the bounding
box $[-T/2, T/2]^4$ and accepts those satisfying $|t| + r \leq T/2$.  The
acceptance rate is $\pi/24 \approx 13\%$ in 4D.  The returned array has exactly
$N$ points in $[t, x, y, z]$ format.

\section{Streaming Chunks}

For the sliding-window path, \func{sprinkle\_chunk()} generates points in a
time slab $[t_{\min}, t_{\max}]$, returning \code{GridPoint} structs with
sequential \code{orig\_idx} numbering.  Grid quantisation fields
(\code{qt, qx, qy, qz}) are filled by the caller after choosing grid
dimensions.


% ███████████████████████████████████████████████████████████████████████████████
\chapter{Phase 1 --- Hasse Diagram}
% ███████████████████████████████████████████████████████████████████████████████

\begin{introduction}
  \item Two-tier Hasse construction: sparse $A^2$ vs.\ geometric OCI
  \item Transitive reduction as the origin of triangle-freeness
  \item Adaptive early termination valves
\end{introduction}

\section{Transitive Reduction}

Source: \file{phase1/hasse.rs}

Given $N$ Poisson-sprinkled events with coordinates $p_i = (t_i, x_i, y_i, z_i)$,
two events are \textbf{causally related} iff
\[
  (t_j - t_i)^2 - |\vec{x}_j - \vec{x}_i|^2 > 0, \qquad t_j > t_i.
\]
The Hasse diagram is the \textbf{transitive reduction}: edge $(i,j)$ exists iff
$i \prec j$ and there is no intermediate $k$ with $i \prec k \prec j$.  This
guarantees triangle-freeness by construction --- any triangle would imply a
transitively redundant edge.

\section{Tier 1: Sparse \texorpdfstring{$A^2$}{A²} Reduction}

For $N \leq 15{,}000$, the function \func{build\_hasse\_sparse()} materialises
the full causal closure as a sparse matrix $A$ (via \code{sprs::TriMat}),
computes $A^2$, and retains only edges where $A^2[i,j] = 0$ (no 2-hop path
exists).  This is exact but $O(N^2)$ in space.

\begin{theorem}[Sparse Reduction Correctness]
  In a transitively closed DAG, edge $(i,j)$ is a Hasse link if and only if
  $A^2[i,j] = 0$.  Proof: if a 2-hop path $i \to k \to j$ exists, then $j$ is
  reachable from $i$ via $k$, making the direct edge redundant.
\end{theorem}

\section{Tier 2: Geometric OCI Construction}

For $N > 15{,}000$, the function \func{build\_hasse\_direct()} avoids
materialising the full closure.  Instead:

\begin{enumerate}
  \item \textbf{Quantise \& sort}: Coordinates are quantised to an integer
        lattice via \func{quantize\_and\_sort()}, yielding \code{GridPoint}
        records sorted by $(q_t, q_x, q_y, q_z)$ for cache locality.
  \item \textbf{Build OCI}: The Occupied Cell Index (OCI) is a per-layer
        structure mapping $(q_x, q_y, q_z)$ cells to the points they contain,
        sorted by $q_x$ for binary search.
  \item \textbf{Parallel link discovery}: For each source $u$, iterate over
        $\Delta t$ shells from 0 to \code{MAX\_CAUSAL\_DEPTH}.  For each
        shell, query the OCI for cells within the light cone
        ($r_{\max} = \Delta t + 1$), test causality, and check transitivity
        against existing children.
  \item \textbf{CSR assembly}: Collect per-chunk degrees and targets, build
        the CSR \code{head} and \code{data} arrays, sort each neighbour list
        for binary-search correctness.
\end{enumerate}

\subsection{Adaptive Termination Valves}

Three early-exit conditions accelerate the search:

\begin{enumerate}
  \item \textbf{Degree cap} (\code{MAX\_HASSE\_DEGREE} $= 15$):
        A node with $\geq 15$ descendants has a fully populated causal shadow.
        Further links are redundant.  This bound also sets the
        \emph{lattice mode decay timescale} for spectral dimension measurements:
        local eigenvalues of the lazy-walk transition matrix near
        $\lambda \approx 1 - 1/(2D)$ decay on a timescale
        $t_{\text{lattice}} = D_{\max}$.
  \item \textbf{Alexandrov volume short-circuit}: If every causal candidate in
        a $\Delta t$ shell is transitively blocked (non-empty Alexandrov set
        between $u$ and $v$), further shells have strictly larger Alexandrov
        volume and will also be blocked.
  \item \textbf{Consecutive dry shells}: Two consecutive shells with no causal
        candidates plus $\geq 2$ existing descendants implies the light cone
        slice is empty.
\end{enumerate}

\section{Bulk Momentum}

Both tiers compute the \textbf{bulk momentum} vector: $p_i = \text{out-degree}(i)
- \text{in-degree}(i)$.  This quantity encodes the net causal asymmetry at each
node and is used in Phase 2 for generation classification ($\phi(w) =
\operatorname{sign}(p_w) \in \{-1, 0, +1\}$).


% ███████████████████████████████████████████████████████████████████████████████
\chapter{Graph Infrastructure}
% ███████████████████████████████████████████████████████████████████████████████

\begin{introduction}
  \item Type-safe CSR with compile-time direction marker
  \item Grid quantisation and the Occupied Cell Index
  \item Physical constants encoded as Rust constants
\end{introduction}

\section{Type-Safe CSR}

Source: \file{graph/csr.rs}

The central data structure is \code{CsrGraph<D>}, a Compressed Sparse Row graph
parametrised by a zero-sized direction marker:

\begin{itemize}
  \item \code{CsrGraph<Directed>}: forward-only edges ($u \to v$ where
        $t_u < t_v$).  Produced by Phase 1.
  \item \code{CsrGraph<Undirected>}: both directions stored explicitly.
        Produced by \func{symmetrize()} or the defect graph.
\end{itemize}

The \code{PhantomData<D>} trick prevents passing a directed Hasse diagram to
spectral walkers (which require the symmetric graph) at \emph{compile time}.

\begin{definition}[CSR Storage Layout]
  \code{head[u]..head[u+1]} indexes into \code{data} to give the sorted
  neighbour list of node $u$.  Binary search on sorted neighbours gives
  $O(\log d)$ edge queries, where $d \leq 15$ in practice.
\end{definition}

Key operations:
\begin{itemize}
  \item \func{neighbors(u)}: returns \code{\&[u32]} slice of sorted neighbours
  \item \func{degree(u)}: returns out-degree (or total degree for undirected)
  \item \func{has\_edge(u, v)}: binary search on sorted neighbour list
  \item \func{symmetrize()}: converts \code{Directed} $\to$ \code{Undirected}
        by adding reverse edges (method on \code{CsrGraph<Directed>} only)
  \item \func{reverse()}: reverses all edge directions (parents $\leftrightarrow$ children)
  \item \func{unique\_edges()}: extracts $(u, v)$ pairs with $u < v$ for
        eigendecomposition (method on \code{CsrGraph<Undirected>} only)
\end{itemize}

Free functions:
\begin{itemize}
  \item \func{build\_directed\_csr(n, rows, cols)}: builds from unsorted edge pairs
  \item \func{build\_undirected\_csr(n, rows, cols)}: builds symmetric CSR from
        edge pairs with automatic reverse insertion
\end{itemize}

\section{Grid Quantisation}

Source: \file{graph/grid.rs}

The OCI construction requires mapping continuous coordinates to an integer lattice.
The \code{GridPoint} struct stores both the original \code{[f64; 4]} coordinates
and quantised \code{(qt, qx, qy, qz)} fields.  The \code{GridDims} struct
encodes lattice dimensions derived from coordinate bounds.

Key constants from \file{grid.rs}:
\begin{itemize}
  \item \code{MAX\_HASSE\_DEGREE = 15}: maximum out-degree in the Hasse diagram.
        Also sets the lattice mode decay timescale $t_{\text{lattice}} = D_{\max}$
        for spectral dimension measurements and the CRT budget floor
        $t_{\max} = D_{\max}$
  \item \code{MAX\_CAUSAL\_DEPTH}: maximum $\Delta t$ shells to search
  \item \code{MAX\_PROPER\_TIME\_SQ}: proper-time cutoff for causality
\end{itemize}


% ███████████████████████████████████████████████████████████████████████████████
\chapter{Phase 2 --- Kuratowski Contraction}
% ███████████████████████████████████████████████████████████████████████████████

\begin{introduction}
  \item Causal Prism ($K_{2,N}$) detection via 2-hop traversal
  \item $K_5$ threat contraction for planarity preservation
  \item Generation classification by bulk-momentum phase signature
  \item Phase-coherence decomposition: visible, dark, sterile mass
\end{introduction}

\section{Causal Prism Detection}

Source: \file{phase2/defect.rs}

The Causal Prism is the fundamental topological particle in FEG.  It is a
complete bipartite subgraph $K_{2,N}$ embedded in the triangle-free Hasse DAG:

\begin{definition}[Causal Prism]
  A \textbf{Causal Prism} consists of:
  \begin{itemize}
    \item An \textbf{origin} $u$ (past pole) and \textbf{destination} $v$ (future pole)
          at graph-distance $\geq 2$
    \item $N \geq 3$ \textbf{intermediaries} $w_1, \ldots, w_N$ such that
          $u \to w_i$ and $w_i \to v$ for all $i$
    \item The intermediaries are mutually disconnected
          (guaranteed by the triangle-free property)
  \end{itemize}
  The minimum $N = 3$ (constant \code{MIN\_PRISM\_SHARED}) derives from
  the Uniqueness of Bifurcation-Convergence theorem: $K_{2,2}$ yields only
  $S_2 \to U(1)$; the $SU(3)$ gauge structure demands $S_3$, hence $N \geq 3$.
\end{definition}

\subsection{Detection Algorithm}

The 2-hop forward-forward traversal:

\begin{enumerate}
  \item Select a \textbf{core} of the top $10\%$ of nodes by undirected degree
        (out-degree + in-degree from the reverse CSR).
  \item For each core node $u$ (not yet placed):
    \begin{enumerate}
      \item Collect 2-hop candidates: for each child $w$ of $u$, for each child
            $v$ of $w$ that is also in the core and not yet placed, add $v$ to
            the candidate set.
      \item For each candidate $v$: compute
            $\text{belly} = \text{children}(u) \cap \text{parents}(v)$ using
            sorted-slice intersection ($O(d_u + d_v)$).
      \item Select the best partner $v^*$ with maximum $|\text{belly}|$.
      \item If $|\text{belly}| \geq 3$: commit the prism, mark all component
            nodes as placed.
    \end{enumerate}
\end{enumerate}

\begin{result}
  The algorithm is $O(N \cdot D^2)$ where $D = $ max Hasse degree $\leq 15$.
  Since $D = O(1)$, the total complexity is $O(N)$.
\end{result}

\section{\texorpdfstring{$K_5$}{K5} Threat Contraction}

After prism detection, external nodes threatening to form a $K_5$ minor are
absorbed:

\begin{definition}[$K_5$ Threat Criterion]
  An external node $t$ is a \textbf{threat} to prism $(u, v, \{w_i\})$ if:
  \begin{itemize}
    \item $t$ is connected (undirected) to \emph{both} poles $u$ and $v$
    \item $t$ is connected to $\geq 2$ intermediaries (constant \code{PRISM\_THREAT})
  \end{itemize}
  The five nodes $\{u, v, t, w_i, w_j\}$ are then mutually reachable, forming a
  $K_5$ minor.  Contraction: $t$ is merged into the higher-degree pole.
\end{definition}

Merge chains are resolved transitively: if $t_1 \to t_2 \to u$, then $t_1$ maps
directly to $u$.  The resulting \code{merge\_map} is a forest of stars.

\section{Generation Classification}

Each intermediate $w_i$ carries a \textbf{causal phase}
$\phi(w_i) = \operatorname{sign}(\text{bulk\_momentum}[w_i]) \in \{-1, 0, +1\}$.

\begin{definition}[Generation Number]
  The \textbf{generation} of a prism $P$ is
  $g(P) = |\{\phi(w_i) : w_i \in \text{intermediaries}\}|$, i.e., the number of
  distinct phase classes.  Since $\phi \in \{-1, 0, +1\}$, $g \in \{1, 2, 3\}$.
\end{definition}

\begin{definition}[Net Phase and Matter/Antimatter]
  The \textbf{net phase} $\Phi(P) = \sum_i \phi(w_i)$.
  \begin{itemize}
    \item $\Phi > 0$: \textbf{matter} (e.g., electron for $g=1$)
    \item $\Phi < 0$: \textbf{antimatter} (e.g., positron for $g=1$)
    \item $\Phi = 0$: \textbf{sterile} (dark matter candidate)
  \end{itemize}
\end{definition}

The generation sets (\code{GenerationSets}) store node indices for Gen1, Gen2,
Gen3, Anti-Gen1, and Sterile prisms.  Average topological mass per generation is
$\langle M \rangle_g = \langle N \rangle$ (average intermediate count).

\section{Phase-Coherence Decomposition}

The defect module computes a decomposition of gravitational mass into visible
and dark components:

\begin{align}
  M_{\text{grav}}(P) &= N, \\
  M_{\text{vis}}(P)  &= |\Phi(P)|, \\
  M_{\text{dark}}(P) &= N - |\Phi(P)|.
\end{align}

Aggregated ratios:
\begin{itemize}
  \item $\Omega_{\text{dark}}/\Omega_{\text{vis}} =
        \sum(N - |\Phi|) / \sum |\Phi|$
  \item $Q_{\text{topo}} = 1/4$ (exact, from topological collider M13)
  \item $\alpha_0 = Q_{\text{topo}} / (8\pi) = 1/(32\pi) \approx 1/100.5$
        (bare Planck coupling; vacuum polarization screens to $\alpha \approx 1/137$)
\end{itemize}

\section{Topology Summary}

Source: \file{phase2/topology.rs}

The \code{TopologySummary} struct aggregates 27 fields including total nodes/prisms,
per-generation counts and masses, phase census, $Q_{\text{topo}}$, $\alpha_{\text{EM}}$,
and the prism histogram (belly-size distribution).


% ███████████████████████████████████████████████████████████████████████████████
\chapter{Phase 3 --- Spectral Observables}
% ███████████████████████████████████████████████████████████████████████████████

\begin{introduction}
  \item Spectral dimension $\dS(t)$ from random walks
  \item Two computation tiers: eigendecomposition vs.\ Monte Carlo walkers
  \item Causal flux: attraction and repulsion between generations
\end{introduction}

\section{Spectral Dimension}

Source: \file{phase3/spectral.rs}

The spectral dimension quantifies the effective dimensionality experienced by a
diffusing particle at scale $t$:
\[
  \dS(t) = -2\,\frac{d \ln P(t)}{d \ln t},
\]
where $P(t)$ is the return probability after $t$ lazy random walk steps on the
graph.  Centred finite differences are used for interior points; forward/backward
differences at endpoints.

The lazy walk (50\% stay, 50\% hop to a random neighbour) has transition matrix
\[
  T = \tfrac{1}{2}(I + D^{-1}A),
\]
where $A$ is the adjacency matrix and $D$ the degree matrix.  Eigenvalues of $T$
lie in $[0,\,1]$ with $\lambda_0 = 1$ for the stationary mode.  The return
probability decomposes as
\[
  P(t) = \frac{1}{N}\operatorname{Tr}(T^t) = \frac{1}{N}\sum_{k} \lambda_k^{t}.
\]
Three eigenvalue bands contribute at different timescales:

\begin{center}
\begin{tabular}{lll}
  \toprule
  \textbf{Band} & \textbf{Eigenvalues} & \textbf{Decay timescale} \\
  \midrule
  Continuum & $\lambda \approx 1$ & $O(N^{2/\dS})$ \\
  Lattice   & $\lambda \approx 1 - O(1/D_{\max})$ & $O(D_{\max}) = O(15)$ \\
  High-freq & $\lambda \approx 0$ & $O(1)$ \\
  \bottomrule
\end{tabular}
\end{center}

After $t \approx D_{\max} = 15$ steps, the lattice modes have decayed by
$\sim 1/e$ and $\dS(t)$ enters the continuum scaling plateau (the
\textbf{lattice-to-continuum crossover}).  This is \emph{not} a mixing time in
the coupon-collector sense (that would be $O(D \ln D) \approx 41$) --- it is the
spectral decay timescale on which discrete graph artifacts (degree echoes,
parity oscillations) damp below the continuum signal.

\begin{definition}[WalkResult]
  Each walk category (vacuum, defect, per-generation, sterile, flux) returns:
  \begin{itemize}
    \item \code{p}: $P(t)$ return probability at each measurement step
    \item \code{ds}: $\dS(t)$ spectral dimension derived from $P(t)$
    \item \code{ds\_std}: ensemble standard deviation (empty for $M = 1$)
  \end{itemize}
\end{definition}

\begin{definition}[SpectralOutput]
  The composed output struct groups:
  \begin{itemize}
    \item Vacuum (global + core) and defect (global + core) walks
    \item Per-generation walks: [Gen1, Gen2, Gen3, Anti1]
    \item Sterile walk (phase-cancelled prisms)
    \item Causal flux: attraction (Gen1$\to$Anti1) and repulsion (Gen1$\to$Gen1)
    \item Normalised flux: flux per target node
    \item Average topological mass per generation
  \end{itemize}
\end{definition}

\section{Tier 1: Eigendecomposition}

For $N \leq$ \code{eigen\_cutoff} (default 3000), the function
\func{compute\_eigen()} computes the full eigendecomposition of the symmetric
normalised adjacency $S = D^{-1/2}AD^{-1/2}$ via \code{nalgebra}.  The
return probability is then
\[
  P(t) = \frac{1}{N}\sum_{k=0}^{N-1} \lambda_k^{t},
\]
where $\lambda_k$ are the eigenvalues of $S$ (clamped to $[-1,1]$).
Core walks restrict the sum to eigenvector-weighted contributions on core nodes.

\section{Tier 2: Monte Carlo Walkers}

Source: \file{phase3/walker.rs}

For $N >$ \code{eigen\_cutoff}, lazy random walkers simulate diffusion.  Each
walker implements the transition matrix $T = \frac{1}{2}(I + D^{-1}A)$
stochastically: at each step it stays with probability $\frac{1}{2}$ or hops to
a uniformly random neighbour with probability $\frac{1}{2}$.  The return
indicator $X_i \in \{0,1\}$ at each measurement step is a $\mathrm{Bernoulli}(P(t))$
variable.  By the CLT, $\mathrm{RSE} = 1/\sqrt{W \cdot P(t)}$.

The number of walkers $W$ is set by the \textbf{Causal Resolution Theorem}:
\[
  W = \lceil (t_{\max} / \varepsilon)^2 \rceil,
\]
where $\varepsilon$ is the topological error tolerance (default 0.01) and
$t_{\max}$ is the maximum diffusion time (default $15 = D_{\max}$).  Setting
$\mathrm{RSE} < \varepsilon$ with $P(t) \sim t^{-2}$ ($\dS = 4$) yields
$W > t^2/\varepsilon^2$.  At $\varepsilon = 0.01$, $t_{\max} = 15$:
$W = 2{,}250{,}000$ walkers.

\textbf{Why $t_{\max} = D_{\max} = 15$:} This is the \emph{lattice mode decay
timescale} --- after $\sim D_{\max}$ steps, the discrete graph artifacts in
$P(t)$ have damped by $\sim 1/e$ and the continuum scaling plateau begins.
Calibrating at the plateau onset is optimal because $P(D_{\max})$ is the
largest return probability in the scaling regime, yielding the smallest required
$W$.  For $t > D_{\max}$, RSE degrades as $(t/D_{\max})\cdot\varepsilon$, but
auto-convergence (Welford RSE $< 5\times 10^{-4}$, $k = 3$ consecutive batches)
adds walker batches until the observable at $t \approx 36$ (midpoint of the step
array, well inside the scaling plateau) stabilises.

Return counts are accumulated as \code{u64} integers (strict finitism).

\section{Causal Flux}

Source: \file{phase3/flux.rs}

The causal flux measures directed transmission between generation sets:
\begin{itemize}
  \item \textbf{Attraction} (Gen1$\to$AntiGen1): flux from matter to antimatter
        (opposite topological charge $\to$ attraction)
  \item \textbf{Repulsion} (Gen1$\to$Gen1): flux from matter to matter
        (same charge $\to$ repulsion)
\end{itemize}
Normalised flux divides by the number of target nodes, isolating the intrinsic
coupling per unit charge.

\section{Diffusion Steps}

The measurement times are densely sampled near $t = 1$:
\begin{verbatim}
  steps = (1..=30).chain((16..=50).map(|i| i * 2))
\end{verbatim}
giving steps 1--30 (unit spacing) and 32, 34, \ldots, 100 (stride 2).  This
covers three regimes:
\begin{itemize}
  \item $t \lesssim D_{\max} = 15$: short-time lattice transient (degree echoes,
        parity oscillations from discrete graph structure)
  \item $D_{\max} \lesssim t \lesssim N^{2/\dS}$: continuum scaling plateau
        where $P(t) \sim t^{-\dS/2}$ and $\dS(t)$ is physically meaningful
  \item $t \gtrsim N^{2/\dS}$: approach to finite-size saturation $P(t) \to 1/N$
\end{itemize}
The convergence observable is $\dS$ at the midpoint of the step array
($t \approx 36$), which lies well within the scaling plateau.


% ███████████████████████████████████████████████████████████████████████████████
\chapter{Ensemble Framework}
% ███████████████████████████████████████████████████████████████████████████████

\begin{introduction}
  \item Adaptive batching with Welford convergence
  \item Checkpointing and resume capability
  \item Seed derivation for deterministic reproducibility
\end{introduction}

\section{Ensemble Runner}

Source: \file{ensemble/runner.rs}

Each ensemble run consists of $M$ independent Monte Carlo realisations of the
full Phase 1--3 pipeline.  The runner uses \textbf{adaptive Welford convergence}:

\begin{enumerate}
  \item Run \code{min\_ensemble} realisations (default 8) unconditionally.
  \item After each subsequent batch, compute the online running mean and
        standard error of \code{mass\_gen1} via the Welford algorithm.
  \item If the relative standard error $\text{SE}/\bar{x} <
        \text{target\_error}$ (default 0.01), declare convergence and stop.
  \item Otherwise, continue up to \code{max\_ensemble} (default 50).
\end{enumerate}

\begin{definition}[EnsembleResult]
  The output contains:
  \begin{itemize}
    \item \code{spectral}: Ensemble-averaged \code{SpectralOutput}
    \item \code{topology}: Aggregated \code{TopologySummary}
    \item \code{measurements}: Aggregated \code{MeasureResults} (if any active)
    \item \code{actual\_m}: Number of realisations actually run
    \item \code{converged}: Whether the target error was reached
  \end{itemize}
\end{definition}

\section{Seed Derivation}

Realisation $k$ uses seed \code{seed\_base + k}.  The base seed is either
user-specified (\flag{seed}) or derived from the system clock.  The
deterministic seeding ensures reproducibility: given the same
\code{(seed, N, M)}, the output is identical on the same platform.

\section{Concurrency}

Realisations within a batch run in parallel via \code{rayon}.  The maximum
concurrency is:
\begin{itemize}
  \item \code{min(batch\_size, max\_concurrent\_runs)}
  \item \code{max\_concurrent\_runs}: for $N > 10^7$, limited to 2; otherwise 4
  \item Overridable via \flag{threads} and \flag{batch-size}
\end{itemize}

\section{Checkpointing}

Source: \file{ensemble/checkpoint.rs}

After each realisation, spectral and topology results are serialised to a
checkpoint file in the output directory.  The \flag{resume} flag reloads the
checkpoint and continues from the last completed realisation.  Provenance
verification prevents cross-fork resume (Layer 4 of the provenance system).

\section{Averaging}

Source: \file{ensemble/averaging.rs}

The \func{average\_ensemble()} function computes element-wise means of
$P(t)$ and $\dS(t)$ vectors across $M$ realisations, plus ensemble standard
deviations.  Topology fields are aggregated by \func{aggregate\_topology()}.


% ███████████████████████████████████████████████████████████████████████████████
\chapter{Measurements M1--M5}
% ███████████████████████████████████████████████████████████████████████████████

\begin{introduction}
  \item M1: Traversal mass ratios
  \item M2: Half-life census
  \item M3: Modulo path integral (NTT)
  \item M4: Vacuum polarisation
  \item M5: Electroweak sector
\end{introduction}

Source: \file{measure/mod.rs}

The measurement system provides fourteen independent physics observables.  Each
measurement is gated by a boolean flag in \code{MeasureFlags} and receives a
\code{MeasureContext} containing the full Phase 1--2 outputs.  Results are
aggregated across ensemble realisations via per-measurement \func{aggregate()}
functions.

\section{M1: Traversal Mass Ratios}

Source: \file{measure/m01\_traversal.rs}

Flag: \flag{measure-mass}

Measures the \textbf{topological inertia} of Causal Prisms by confining random
walkers to prism components and recording cover times.  The ratio of cover times
between generations gives the mass hierarchy:
\[
  \frac{m_\mu}{m_e} \approx \frac{\langle T_{\text{cover}} \rangle_{\text{Gen2}}}
  {\langle T_{\text{cover}} \rangle_{\text{Gen1}}}.
\]

Output: \file{traversal\_mass.csv} with per-generation cover times and ratios.

\section{M2: Half-Life Census}

Source: \file{measure/m02\_halflife.rs}

Flag: \flag{measure-halflife}

Measures the topological stability of prisms by recording edge-snap events
--- the random removal of one intermediary edge that splits a $K_{2,N}$ into
$K_{2,N-1} + K_{2,1}$.  The half-life is the Markov hitting time for the first
snap event at a given thermal excitation.

\begin{result}
  The electron ($K_{2,3}$) is topologically stable: splitting to $K_{2,2} +
  K_{2,1}$ requires overcoming the $K_5$ planarity barrier (energy $\sim M_{\text{Pl}}$).
  The muon and tau have finite half-lives proportional to their genus gap.
\end{result}

Output: \file{half\_life.csv} with per-generation half-lives and snap rates.

\section{M3: Modulo Path Integral (NTT)}

Source: \file{measure/m03\_modulo.rs}

Flag: \flag{measure-modulo}

Implements the \textbf{Number-Theoretic Transform} (NTT) path integral: walkers
accumulate action $S$ as $g^S \bmod p$ where $g$ is a primitive root and $p$ is
a prime modulus.  This replaces the continuum $e^{iS/\hbar}$ with a discrete
$\mathbb{F}_p$ phase.

Configuration:
\begin{itemize}
  \item \code{prime}: default $65537 = 2^{16} + 1$ (Fermat prime)
  \item \code{root}: default $3$ (primitive root of $65537$)
  \item \code{steps}: default $500$ NTT walk steps
\end{itemize}

The interference pattern emerges from modular cancellation: when $\sum g^{S_j}
\equiv 0 \pmod{p}$, the path is destructively forbidden.

Output: \file{modulo\_interference.csv} with phase histograms and interference fringe
visibility.

\section{M4: Vacuum Polarisation}

Source: \file{measure/m04\_vacuum\_pol.rs}

Flag: \flag{measure-vacuum}

Measures the running of the topological charge $Q_{\text{topo}}$ with scale.
At short distances (small walker steps), the effective coupling increases
due to virtual prism--anti-prism screening:
\[
  \alpha_{\text{EM}}^{-1}(0) \approx 132 \quad \to \quad
  \alpha_{\text{EM}}^{-1}(m_Z) \approx 128.
\]

The mechanism is planarity screening: at short distances, prism--anti-prism pairs
are resolved and screen the effective charge.

Output: \file{vacuum\_polarization.csv} with $Q_{\text{topo}}(t)$ and
$\alpha^{-1}(t)$ at each scale $t$.

\section{M5: Electroweak Sector}

Source: \file{measure/m05\_electroweak.rs}

Flag: \flag{measure-electroweak}

Probes the separation of $SU(2)_L \times U(1)_Y$ from the prism topology.
This measurement uses the asymmetry between left-handed and right-handed walker
trajectories to extract the Weinberg angle $\theta_W$:
\[
  \sin^2\theta_W = \frac{\text{left-blocked}}{\text{left-blocked} +
  \text{right-blocked}}.
\]

Output: \file{electroweak.csv} with mixing angle and left/right asymmetry.


% ███████████████████████████████████████████████████████████████████████████████
\chapter{Measurements M6--M10}
% ███████████████████████████████████████████████████████████████████████████████

\begin{introduction}
  \item M6: Quantum decoherence and Born rule
  \item M7: Neutrino census
  \item M8: PMNS mixing matrix
  \item M9: Higgs mechanism (topological drag)
  \item M10: Full Standard Model Lagrangian card
\end{introduction}

\section{M6: Quantum Decoherence}

Source: \file{measure/m06\_decoherence.rs}

Flag: \flag{measure-decoherence}

Tests the emergence of the Born rule via the Cramer transactional handshake.
A retarded wave (fwd) is seeded from prism intermediates and propagated forward
with NTT phase rotation $W = g^{(p-1)/4} \bmod p$ (90° per edge); an advanced
wave (bwd) is seeded from the future boundary and propagated backward with
$W^{-1}$.  The handshake product, centered symmetrically in $\mathbb{Z}/p\mathbb{Z}$,
gives the Born amplitude at each environment node:
\[
  P(d) \;\propto\; \bigl|\,\text{fwd}(d) \times \text{bwd}(d)\,\bigr|^2.
\]
$K_{2,2}$ diamond counts at environment nodes provide the observed detection
events.  Pearson~$r$ between the predicted $|\psi|^2$ PMF and the observed
$K_{2,2}$ PMF, validated against a permutation null model (1000~shuffles),
measures Born rule compliance.

Output: \file{decoherence.csv} and \file{born\_rule.csv} with Pearson~$r$,
null model statistics, and per-bin predicted vs.\ observed frequencies.

\section{M7: Neutrino Census}

Source: \file{measure/m07\_neutrino.rs}

Flag: \flag{measure-neutrino}

Counts ``open prisms'' --- $K_{2,1}$ structures (two poles, one intermediary)
that are topologically sterile ($\Phi = 0$) and have near-zero topological mass.
These are identified as neutrinos: gravitationally active but electromagnetically
invisible.

The census provides per-generation neutrino counts and the
neutrino/charged-lepton ratio.

Output: \file{neutrino.csv} with generation-resolved neutrino counts.

\section{M8: PMNS Mixing Matrix}

Source: \file{measure/m08\_pmns.rs}

Flag: \flag{measure-pmns}

Requires M7 (automatically enabled).  Extracts the PMNS mixing matrix from
neutrino oscillation probabilities measured via walker transmission between
generation-tagged open prisms.  The mixing angles $(\theta_{12}, \theta_{13},
\theta_{23})$ are extracted from the matrix elements.

\begin{result}
  FEG predicts ``anarchy''-style PMNS mixing: large mixing angles emerging from
  the random topology of the Hasse diagram, in agreement with the observed
  near-maximal $\theta_{23}$ and large $\theta_{12}$.
\end{result}

Output: \file{pmns.csv} with the $3 \times 3$ mixing matrix and extracted angles.

\section{M9: Higgs Mechanism}

Source: \file{measure/m09\_higgs.rs}

Flag: \flag{measure-higgs}

Models the Higgs boson as a ``breathing mode'' of the Causal Prism ensemble:
a coherent oscillation of the average belly size $\langle N \rangle$ about its
equilibrium value.  The entropy cost $\Delta S \sim 1/\alpha$ of a uniform shift
in $N$ sets the Higgs mass scale:
\[
  m_H \sim v \sqrt{2\lambda}, \qquad \lambda \sim \alpha^2.
\]

The measurement applies topological drag --- walkers on prism components experience
an effective potential proportional to their belly size, and the oscillation
frequency is extracted from the autocorrelation.

Output: \file{higgs.csv} with breathing-mode frequency, effective $\lambda$,
and mass estimate.

\section{M10: Standard Model Lagrangian Card}

Source: \file{measure/m10\_lagrangian.rs}

Flag: \flag{measure-lagrangian}

Assembles a complete Standard Model parameter card from M1--M9 outputs:

\begin{longtable}{lll}
  \toprule
  \textbf{Parameter} & \textbf{Source} & \textbf{FEG Derivation} \\
  \midrule
  $\alpha_0 = 1/(32\pi)$ & M13 & Topological collider $Q_{\text{topo}} = 1/4$ \\
  $\sin^2\theta_W$     & M5 & Left/right walker asymmetry \\
  $m_e : m_\mu : m_\tau$ & M1 & Cover-time ratios \\
  PMNS matrix           & M8 & Neutrino walker transmission \\
  $m_H$                 & M9 & Breathing-mode frequency \\
  $\tau_\mu, \tau_\tau$  & M2 & Edge-snap hitting times \\
  $\alpha_s$            & M4 & Short-distance $Q_{\text{topo}}$ running \\
  $P = |K|^2$           & M6 & Modular path integral \\
  Neutrino counts       & M7 & Open-prism census \\
  \bottomrule
\end{longtable}

This is a post-hoc assembly: M10 does not run its own simulation, but collects
and cross-validates the outputs of M1--M9.

Output: \file{lagrangian.csv} with the full parameter card.


% ███████████████████████████████████████████████████████████████████████████████
\chapter{Measurements M11--M14}
% ███████████████████████████████████████████████████████████████████████████████

\begin{introduction}
  \item M11: Hausdorff dimension via BFS volume growth
  \item M12: Zigzag Kaluza--Klein dimension
  \item M13: Topological particle collider ($Q_{\text{topo}} = 1/4$)
  \item M14: Exact mass formula (genus ladder)
\end{introduction}

These four measurements extract deterministic topological observables that
do not require Phase 3 random walks.  They are automatically enabled by
\flag{topology-only} mode and are also included in \flag{measure-all}.

\section{M11: Hausdorff Dimension}

Source: \file{measure/m11\_hausdorff.rs}

Flag: \flag{measure-hausdorff}

Measures the Hausdorff dimension $d_H$ by sampling random source nodes and
performing BFS on both the directed (forward-only) and undirected (symmetric)
vacuum CSR.  Volume growth $V(r) \sim r^{d_H}$ is fitted via log-log linear
regression over the power-law region ($r = 1 \ldots 8$, before saturation):
\[
  d_H = \frac{d \ln V}{d \ln r} \approx 4.0 \quad \text{(directed BFS)}.
\]

The spectral dimension $\dS \approx 2$ at UV is a lazy-walk artifact on the
tree-like Hasse DAG; Hausdorff dimension via volume growth is the correct
geometric observable for the causal set.

\begin{result}
  At $N = 10^7$, $M = 1$: $d_H(\text{directed}) = 4.0$ (single realisation;
  ensemble average converges to 4.0 exactly).
\end{result}

Output: \file{hausdorff.csv} with $V(r)$ at each BFS depth and fitted $d_H$.

\section{M12: Zigzag KK Dimension}

Source: \file{measure/m12\_zigzag.rs}

Flag: \flag{measure-zigzag}

Measures the zigzag Kaluza--Klein dimension $d_{\text{zig}}$ by alternating
forward and backward BFS steps on the directed vacuum CSR and its reverse.
The total volume growth $V(r)$ is fitted as $V \sim r^{d_{\text{zig}}}$.

The zigzag path explores both future and past causal directions, effectively
probing the dimensionality of the full causal structure including KK-like
extra directions:
\[
  d_{\text{zig}} \approx 4.004.
\]

The screening ratio $S_4 / S_{d_{\text{zig}}} = 32\pi / 137 \approx 0.733$
connects the bare coupling $\alpha_0 = 1/(32\pi)$ to the observed
$\alpha \approx 1/137$.

Output: \file{zigzag.csv} with zigzag $V(r)$ and fitted $d_{\text{zig}}$.

\section{M13: Topological Particle Collider}

Source: \file{measure/m13\_collider.rs}

Flag: \flag{measure-collider}

Extracts $Q_{\text{topo}}$ from a strict bipartite seam plus external Markov
blanket analysis of $K_{2,N}$ prisms.  For each committed prism:

\begin{enumerate}
  \item \textbf{Seam}: the exclusive children of each pole, filtered by
        $z > \text{emitting\_top}$ (no tachyons)
  \item \textbf{Blanket}: edges from diamond nodes to non-diamond nodes only
        (strict external boundary)
  \item $Q_{\text{topo}} = \text{seam\_size} / \text{blanket\_size}$
\end{enumerate}

\begin{result}
  $Q_{\text{topo}} = 1/4$ (exact, locked asymptote across $N = 10^4$ to $10^7$).
  The bare coupling is $\alpha_0 = Q_{\text{topo}} / (8\pi) = 1/(32\pi)
  \approx 1/100.5$.  This supersedes the earlier ``port counting'' method
  ($Q_{\text{topo}} = 4/21$) used in earlier versions.
\end{result}

Output: \file{collider.csv} with per-prism $Q_{\text{topo}}$ and ensemble mean.

\section{M14: Exact Mass Formula}

Source: \file{measure/m14\_mass\_formula.rs}

Flag: \flag{measure-mass-formula}

Computes the genus $g$ of every $K_{2,n}$ prism via the Grothendieck--Euler
formula (\file{phase2/writhe.rs}) and verifies the exact mass formula:
\[
  M(g) = M_e \cdot 2^g \cdot (32\pi)^g \;\big/\; \binom{\max(3,\,2g+1)}{3}.
\]

Three factors:
\begin{itemize}
  \item $2^g$: parity filter (binary Kuratowski junction per hole)
  \item $(32\pi)^g = (1/\alpha_0)^g$: parallel homology drag (metric wrap per hole)
  \item $\binom{n_{\min}}{3}$: subgraph superposition ($K_{2,3}$ backbone count
        in $K_{2,n}$), where $n_{\min}$ from the Grothendieck--Euler floor
        $F = n - 2g \geq 1$
\end{itemize}

\begin{result}
  Zero free parameters:
  \begin{itemize}
    \item $m_\mu / m_e = 201.06$ \quad (SM: 206.77, $-2.76\%$)
    \item $m_\tau / m_e = 4042.6$ \quad (SM: 3477.5, $+16.2\%$)
  \end{itemize}
  Genus classification: $g = 0 \to$ electron, $g = 1 \to$ muon, $g = 2 \to$ tau.
\end{result}

Output: \file{mass\_formula.csv} with genus histogram, mean belly per genus,
and predicted mass ratios.


% ███████████████████████████████████████████████████████████████████████████████
\chapter{Output \& CSV}
% ███████████████████████████████████████████████████████████████████████████████

\begin{introduction}
  \item CSV serialisation with provenance headers
  \item 36-column spectral CSV format
  \item Topology summary and mass spectrum CSVs
  \item Human-readable terminal summary
\end{introduction}

\section{CSV Writer}

Source: \file{output/csv.rs}

The \code{CsvWriter} struct wraps a \code{BufWriter<File>} and prepends the
standard provenance header block (Layer 3):

\begin{verbatim}
  # FEG Kuratowski Calculus Engine (FEG_prism)
  # Author: Juan Pablo Silva Alvarado
  # DOI: 10.5281/zenodo.18769707
  # Framework: Fractal Entropic Geometrodynamics
  # Provenance SHA-256: a4fee337...
  # Commit: 5b07174 (dirty)
  # Timestamp: 2026-02-24T12:00:00Z
\end{verbatim}

\section{Spectral CSV}

Source: \file{output/mod.rs}

The function \func{write\_spectral\_csv()} produces \file{results.csv} with
36 columns per diffusion step:

\begin{longtable}{lp{9cm}}
  \toprule
  \textbf{Column(s)} & \textbf{Description} \\
  \midrule
  \code{step}           & Diffusion time $t$ (integer steps) \\
  \code{P\_vac, dS\_vac}  & Vacuum return probability and spectral dimension \\
  \code{P\_def, dS\_def}  & Defect (post-contraction) return prob.\ and $\dS$ \\
  \code{P\_Gen1..3, dS\_Gen1..3} & Per-generation walks \\
  \code{P\_Anti1, dS\_Anti1}     & Antimatter generation walk \\
  \code{Flux\_Attr, Flux\_Repu}  & Causal flux (raw) \\
  \code{Flux\_Attr\_Norm, Flux\_Repu\_Norm} & Normalised flux per target \\
  \code{P\_Sterile, dS\_Sterile} & Sterile prism walk \\
  \code{Mass\_Gen1..3, Mass\_Anti1} & Average topological mass \\
  \code{dS\_*\_std}       & Ensemble standard deviations \\
  \code{Flux\_*\_std}     & Flux ensemble standard deviations \\
  \bottomrule
\end{longtable}

\section{Topology CSV}

The function \func{write\_topology\_csv()} writes \file{topology\_summary.csv}
with all 27 \code{TopologySummary} fields.  The function
\func{write\_mass\_spectrum\_csv()} writes \file{mass\_spectrum.csv} with the
prism histogram (belly size $\to$ count).

\section{Measurement CSVs}

The function \func{measure::write\_all\_csv()} dispatches to per-measurement
\func{write\_csv()} functions.  Output files:

\begin{itemize}
  \item \file{traversal\_mass.csv} (M1)
  \item \file{half\_life.csv} (M2)
  \item \file{modulo\_interference.csv} (M3)
  \item \file{vacuum\_polarization.csv} (M4)
  \item \file{electroweak.csv} (M5)
  \item \file{decoherence.csv}, \file{born\_rule.csv} (M6)
  \item \file{neutrino.csv} (M7)
  \item \file{pmns.csv} (M8)
  \item \file{higgs.csv} (M9)
  \item \file{lagrangian.csv} (M10)
  \item \file{hausdorff.csv} (M11)
  \item \file{zigzag.csv} (M12)
  \item \file{collider.csv} (M13)
  \item \file{mass\_formula.csv} (M14)
\end{itemize}

\section{Terminal Summary}

Source: \file{output/summary.rs}

The \func{print\_summary()} function writes a human-readable report to stdout
including: spectral dimension at key scales, mass hierarchy ratios, prism counts,
$\alpha_{\text{EM}}$, dark matter fraction, and measurement highlights.

\section{Metadata}

The \code{Metadata} struct accompanies every CSV file:
\begin{itemize}
  \item \code{n\_points}, \code{actual\_m}, \code{converged}
  \item \code{min\_ensemble}, \code{max\_ensemble}
  \item \code{mode} (``in-memory'' or ``streaming'')
  \item \code{epsilon}, \code{tmax}, \code{walkers}
  \item \code{seed}, \code{timestamp}, \code{commit}
\end{itemize}


% ███████████████████████████████████████████████████████████████████████████████
\chapter{Animation Export}
% ███████████████████████████████████████████████████████████████████████████████

\begin{introduction}
  \item Binary topology slice for Manim scenes
  \item Compact bincode serialisation
  \item Integration with the animation layer
\end{introduction}

\section{Topology Slice}

Source: \file{anim\_export.rs}

The animation export produces a single-realisation snapshot containing:

\begin{definition}[TopologySlice]
  \begin{itemize}
    \item \code{n\_total}: number of spacetime events
    \item \code{coordinates}: $[t, x, y, z]$ for each event
    \item \code{hasse\_edges}: directed $(u, v)$ Hasse links
    \item \code{prisms}: vector of \code{PrismDef} (origin, destination,
          intermediates)
  \end{itemize}
\end{definition}

The slice is serialised via \code{bincode} to a compact binary file.  Typical
size: $\sim 1\,\text{MB}$ for $N = 5000$.

\section{Usage}

Invoked by the \flag{export-slice} flag:

\begin{verbatim}
  feg_prism 5000 1 ./output --inmemory \
    --export-slice ./output/slice.bin --seed 42
\end{verbatim}

This runs a single realisation (Phases 1--2 only, no spectral walkers) and
writes the topology slice.  The Python animation layer
(\file{animations/causal\_anim/scene.py}) loads the slice and renders:

\begin{itemize}
  \item 3D Hasse diagram with prism highlighting
  \item Time-evolution of causal structure
  \item Generation-coloured prism components
\end{itemize}


% ███████████████████████████████████████████████████████████████████████████████
\chapter{Configuration}
% ███████████████████████████████████████████████████████████████████████████████

\begin{introduction}
  \item TOML configuration file \file{causal\_set.toml}
  \item RAM-aware mode selection
  \item Memory estimation formulas
\end{introduction}

\section{Configuration File}

Source: \file{config.rs}

The engine reads \file{causal\_set.toml} from the output directory or current
directory.  If absent, a default is created:

\begin{verbatim}
  # FEG Prism Simulation — configuration
  output_dir = "."
  max_ram_gb = 0        # 0 = auto-detect
  safety_fraction = 0.70
\end{verbatim}

\section{Config Struct}

\begin{definition}[Config Fields]
  \begin{itemize}
    \item \code{max\_ram\_bytes}: hard RAM limit (0 = use system detection)
    \item \code{safety\_fraction}: fraction of available RAM considered safe
          (default 0.70, clamped to $[0.05, 1.0]$)
    \item \code{output\_dir}: default output directory (CLI overrides)
  \end{itemize}
\end{definition}

\section{Memory Estimation}

\begin{definition}[Memory Formulas]
  In-memory mode:
  \[
    \text{est}(N) = \begin{cases}
      N^2 \times 24  & N \leq 3000 \text{ (eigendecomp: dense matrix)} \\
      N \times 2000  & 3000 < N \leq 50{,}000 \\
      N \times 500   & N > 50{,}000 \text{ (Hasse degree $\leq 15$)}
    \end{cases}
  \]
  Streaming mode:
  \[
    \text{est}(N) = N \times 13 + 100\,\text{MB}
  \]
\end{definition}

The recommendation engine (\func{recommend\_mode()}) computes
$\text{est}(N) \times \text{concurrent}$ and compares to
$\text{available RAM} \times \text{safety\_fraction}$.

\section{Mode Selection Flow}

\begin{enumerate}
  \item If \flag{inmemory}: force in-memory.
  \item If \flag{stream}: force streaming.
  \item Otherwise: auto-detect via \func{recommend\_mode()}.  Display memory
        summary and prompt user to confirm or override.
\end{enumerate}


% ███████████████████████████████████████████████████████████████████████████████
\chapter{Provenance}
% ███████████████████████████████████████████████████████████████████████████████

\begin{introduction}
  \item Five-layer SHA-256 provenance system
  \item No external crate for SHA-256 (custom implementation)
  \item Binding the binary to its author, DOI, and git commit
\end{introduction}

\section{Architecture}

Source: \file{provenance.rs}

The provenance system provides a chain of trust from source code to output files:

\begin{longtable}{clp{8cm}}
  \toprule
  \textbf{Layer} & \textbf{When} & \textbf{Purpose} \\
  \midrule
  1 & Compile & \code{\#[test]} canary: \code{sha256(PREIMAGE) == HASH} \\
  2 & Startup & Runtime verification: \func{verify\_provenance()} \\
  3 & Output  & CSV header stamps with author, DOI, hash, commit \\
  4 & Resume  & Checkpoint provenance: reject cross-fork resume \\
  5 & Build   & \code{env!()} embedding: git hash + dirty flag \\
  \bottomrule
\end{longtable}

\section{Preimage and Hash}

The provenance preimage is the ASCII string:

\begin{verbatim}
  Juan Pablo Silva Alvarado:10.5281/zenodo.18733424:
  FEG-Kuratowski-2026
\end{verbatim}

Its SHA-256 hash is:

\begin{verbatim}
  a4fee33737305b5eb78a16be2ff233bd
  ce125bffa352127fcbcd997d018c05aa
\end{verbatim}

Both the preimage and hash are compiled into the binary as constants.

\section{Custom SHA-256}

The engine implements SHA-256 from scratch without external crates (204 lines).
The implementation:
\begin{itemize}
  \item Supports messages up to 119 bytes (sufficient for the preimage)
  \item Follows FIPS 180-4 exactly: padding, schedule expansion, 64-round
        compression with the standard $K$ constants
  \item Verified by the compile-time canary and two additional unit tests
        (\func{hex\_matches\_bytes}, \func{runtime\_verify})
\end{itemize}

\section{Build-Time Embedding}

The \code{build.rs} script (not shown) sets two environment variables:
\begin{itemize}
  \item \code{GIT\_HASH}: short git commit hash
  \item \code{GIT\_DIRTY}: ``dirty'' or ``clean''
\end{itemize}
These are embedded via \code{env!()} macros and surfaced by
\func{commit\_string()}: e.g., \code{"5b07174 (dirty)"}.

\section{UTC Timestamp}

The function \func{utc\_timestamp()} computes the current time without external
crates, using \code{SystemTime::now()} and manual calendar arithmetic
(leap-year-aware).  Format: ISO 8601 (\code{2026-02-24T12:00:00Z}).


% ███████████████████████████████████████████████████████████████████████████████
\chapter{Usage \& CLI Reference}
% ███████████████████████████████████████████████████████████████████████████████

\begin{introduction}
  \item Complete CLI reference with all flags
  \item Quick-start examples
  \item Common workflows
\end{introduction}

\section{Synopsis}

\begin{verbatim}
  feg_prism [N] [M] [output_dir] [FLAGS...]
\end{verbatim}

\subsection{Positional Arguments}

\begin{longtable}{llp{8cm}}
  \toprule
  \textbf{Argument} & \textbf{Default} & \textbf{Description} \\
  \midrule
  \code{N}          & 5000   & Number of spacetime events \\
  \code{M}          & 50     & Maximum ensemble realisations (or use \flag{max-ensemble}) \\
  \code{output\_dir} & \code{.} & Directory for all output files \\
  \bottomrule
\end{longtable}

\section{Flags}

\subsection{Execution Mode}

\begin{longtable}{lp{10cm}}
  \toprule
  \textbf{Flag} & \textbf{Description} \\
  \midrule
  \flag{inmemory}  & Force in-memory mode (fast) \\
  \flag{stream}    & Force streaming mode (low RAM) \\
  \flag{resume}    & Resume from last checkpoint \\
  \bottomrule
\end{longtable}

\subsection{Physics Parameters}

\begin{longtable}{llp{8cm}}
  \toprule
  \textbf{Flag} & \textbf{Default} & \textbf{Description} \\
  \midrule
  \flag{epsilon} \code{<f64>}  & 0.01  & Topological error tolerance $\varepsilon$ \\
  \flag{tmax} \code{<usize>}   & 15    & CRT budget calibration time ($= D_{\max}$, lattice mode decay timescale) \\
  \flag{seed} \code{<u64>}     & clock & Base seed (deterministic mode) \\
  \flag{eigen-cutoff} \code{<N>} & 3000 & Eigendecomp threshold \\
  \bottomrule
\end{longtable}

\subsection{Ensemble Control}

\begin{longtable}{llp{8cm}}
  \toprule
  \textbf{Flag} & \textbf{Default} & \textbf{Description} \\
  \midrule
  \flag{max-ensemble} \code{<N>} & 50   & Safety cap on realisations \\
  \flag{min-ensemble} \code{<N>} & 8    & Minimum before convergence check \\
  \flag{target-error} \code{<F>} & 0.01 & Relative SE target \\
  \flag{batch-size} \code{<N>}   & 4    & Max concurrent per batch \\
  \flag{threads} \code{<N>}      & auto & Max parallel realisations \\
  \bottomrule
\end{longtable}

\subsection{Measurements}

\begin{longtable}{lp{10cm}}
  \toprule
  \textbf{Flag} & \textbf{Measurement} \\
  \midrule
  \flag{measure-mass}        & M1: Traversal mass ratios \\
  \flag{measure-halflife}    & M2: Half-life census \\
  \flag{measure-modulo}      & M3: Modulo path integral \\
  \flag{measure-vacuum}      & M4: Vacuum polarisation \\
  \flag{measure-electroweak} & M5: Electroweak sector \\
  \flag{measure-decoherence} & M6: Quantum decoherence \\
  \flag{measure-neutrino}    & M7: Neutrino census \\
  \flag{measure-pmns}        & M8: PMNS mixing matrix (enables M7) \\
  \flag{measure-higgs}       & M9: Higgs mechanism \\
  \flag{measure-lagrangian}  & M10: Full SM Lagrangian card (enables M1--M9) \\
  \flag{measure-hausdorff}   & M11: Hausdorff dimension (BFS volume growth) \\
  \flag{measure-zigzag}      & M12: Zigzag KK dimension \\
  \flag{measure-collider}    & M13: Topological collider ($Q_{\text{topo}}$) \\
  \flag{measure-mass-formula} & M14: Exact mass formula (genus ladder) \\
  \flag{measure-all}         & Enable all M1--M9, M11--M14 \\
  \flag{topology-only}       & Skip Phase 3 walks; auto-enable M11/M13/M14 ($\sim$5$\times$ faster) \\
  \bottomrule
\end{longtable}

\subsection{Modulo Configuration}

\begin{longtable}{llp{8cm}}
  \toprule
  \textbf{Flag} & \textbf{Default} & \textbf{Description} \\
  \midrule
  \flag{modulo-prime} \code{<u64>} & 65537 & Prime modulus $p$ for M3/M6 \\
  \flag{modulo-root} \code{<u64>}  & 3     & Primitive root $g$ of $p$ \\
  \flag{modulo-steps} \code{<N>}   & 500   & NTT walk steps \\
  \bottomrule
\end{longtable}

\subsection{Export}

\begin{longtable}{lp{10cm}}
  \toprule
  \textbf{Flag} & \textbf{Description} \\
  \midrule
  \flag{export-slice} \code{<PATH>} & Export topology slice (single realisation, early exit) \\
  \flag{help}                        & Show usage help \\
  \bottomrule
\end{longtable}

\section{Quick-Start Examples}

\subsection{Minimal Run}

\begin{verbatim}
  cargo run --release --bin feg_prism -- \
    5000 8 ./output --inmemory --seed 42
\end{verbatim}

Runs $N = 5000$, $M \leq 8$ realisations, deterministic seed, in-memory mode.

\subsection{Full Standard Model Card}

\begin{verbatim}
  cargo run --release --bin feg_prism -- \
    50000 25 ./sm_run --inmemory \
    --measure-lagrangian --seed 42
\end{verbatim}

Enables all M1--M10 measurements and produces the complete SM parameter card.

\subsection{Production Run}

\begin{verbatim}
  cargo run --release --bin feg_prism -- \
    1000000 50 ./production --inmemory \
    --measure-all --target-error 0.005 --seed 42
\end{verbatim}

One million events, up to 50 realisations, tighter convergence target.

\subsection{Born Rule Verification}

\begin{verbatim}
  cargo run --release --bin feg_prism -- \
    5000 10 ./born --inmemory \
    --measure-decoherence --seed 42
\end{verbatim}

\subsection{PMNS Mixing Matrix}

\begin{verbatim}
  cargo run --release --bin feg_prism -- \
    50000 25 ./pmns --inmemory \
    --measure-pmns --seed 42
\end{verbatim}

Automatically enables M7 (neutrino census) as a prerequisite.

\subsection{Higgs Mass Estimate}

\begin{verbatim}
  cargo run --release --bin feg_prism -- \
    50000 25 ./higgs --inmemory \
    --measure-higgs --seed 42
\end{verbatim}

\subsection{Animation Export}

\begin{verbatim}
  cargo run --release --bin feg_prism -- \
    5000 1 ./anim --inmemory \
    --export-slice ./anim/slice.bin --seed 42
\end{verbatim}

Early exit after Phase 2; no spectral walkers run.

\subsection{Resume from Checkpoint}

\begin{verbatim}
  cargo run --release --bin feg_prism -- \
    50000 50 ./production --inmemory \
    --measure-all --resume --seed 42
\end{verbatim}

Continues from the last checkpoint in \file{./production/}.

\subsection{Custom Modulo Parameters}

\begin{verbatim}
  cargo run --release --bin feg_prism -- \
    5000 10 ./modulo --inmemory \
    --measure-modulo --modulo-prime 257 \
    --modulo-root 3 --modulo-steps 1000 --seed 42
\end{verbatim}


% ███████████████████████████████████████████████████████████████████████████████
\chapter{Expected Results \& Verification}
% ███████████████████████████████████████████████████████████████████████████████

\begin{introduction}
  \item Expected output values at standard scales
  \item Verification checklist
  \item Known finite-size effects
\end{introduction}

\section{Spectral Dimension}

\begin{result}
  At $N = 50{,}000$, $M = 25$:
  \begin{itemize}
    \item $\dS(t \to 0) \to 2$ (short-distance fractal)
    \item $\dS(t \to \infty) \to 4$ (long-distance manifold recovery)
    \item The lattice-to-continuum crossover occurs at $t \sim D_{\max} = 15$
          diffusion steps, consistent with the lattice mode decay timescale.
          Below this, $\dS(t)$ exhibits lattice artifacts (degree echoes,
          parity oscillations); above this, $P(t) \sim t^{-2}$ scaling emerges
    \item Defect $\dS$ is slightly reduced compared to vacuum
          (matter introduces topological shortcuts)
  \end{itemize}
\end{result}

\section{Mass Hierarchy}

\begin{result}
  Cover-time mass ratios (M1, $N = 50{,}000$, $M = 25$):
  \begin{itemize}
    \item $m_\mu / m_e \approx 200 \pm 30$
          (target: 206.77, PDG)
    \item $m_\tau / m_e \approx 3500 \pm 500$
          (target: 3477.4, PDG)
  \end{itemize}
  Exact agreement improves with $N$; finite-size corrections scale as $N^{-1/4}$.
\end{result}

\section{Fine Structure Constant}

\begin{result}
  Topological collider (strict bipartite seam + external Markov blanket):
  \begin{itemize}
    \item $Q_{\text{topo}} = 1/4$ (exact, locked asymptote across $N = 10^4$ to $10^7$)
    \item $\alpha_0 = Q_{\text{topo}} / (8\pi) = 1/(32\pi) \approx 1/100.5$ (bare Planck coupling)
    \item Screening ratio $32\pi/137 \approx 0.733$: vacuum polarization to lab value
  \end{itemize}
\end{result}

\section{Dark Matter Fraction}

\begin{result}
  Phase-coherence decomposition ($N = 50{,}000$, $M = 25$):
  \begin{itemize}
    \item $\Omega_{\text{dark}} / \Omega_{\text{vis}} \approx 4$--$6$
          (target: 5.3, Planck)
    \item Sterile prisms ($\Phi = 0$) constitute $\sim 10$--$20\%$ of all prisms
  \end{itemize}
\end{result}

\section{Born Rule}

\begin{result}
  Decoherence measurement (M6, $N = 5{,}000$, $M = 10$):
  \begin{itemize}
    \item Pearson~$r$ between predicted $|\psi|^2$ and observed $K_{2,2}$ PMF validates Born rule
    \item Permutation null model (1000 shuffles) confirms significance
    \item Coherence decay correlation vs $1/\sqrt{\text{env\_size}}$ measures decoherence rate
  \end{itemize}
\end{result}

\section{Verification Checklist}

\begin{enumerate}
  \item Compile: \code{cargo build --release --bin feg\_prism}
  \item Provenance: startup message shows ``SHA-256 verified $\checkmark$''
  \item Triangle-freeness: \code{topology\_summary.csv} shows zero $K_4$ count
  \item Prism counts: $\sim 5$--$15\%$ of core nodes form prisms
  \item Generation hierarchy: Gen1 > Gen2 > Gen3 in abundance
  \item Spectral dimension: $\dS(1) \approx 2$, $\dS(30) \approx 4$
  \item CSV files: all provenance headers present and hash matches
  \item Determinism: same seed produces identical output
\end{enumerate}

\section{Known Finite-Size Effects}

\begin{itemize}
  \item \textbf{Boundary effects}: The causal diamond has sharp boundaries.
        Nodes near the tips have anomalously low degree.  The core selection
        (top 10\% by degree) mitigates this by preferring interior nodes.
  \item \textbf{Mass ratio convergence}: The cover-time ratios $m_\mu/m_e$ and
        $m_\tau/m_e$ converge slowly (as $N^{-1/4}$).  For precision to within
        $10\%$, use $N \geq 50{,}000$.
  \item \textbf{Spectral dimension plateau}: The $\dS \to 4$ plateau requires
        $t \gtrsim D_{\max} = 15$ (lattice mode decay timescale).  The default
        $t_{\max} = 15$ calibrates the CRT budget at the onset of this plateau;
        the full step schedule extends to $t = 100$ and auto-convergence ensures
        the observable at $t \approx 36$ stabilises.
  \item \textbf{Ensemble variance}: Some observables (e.g., $\alpha_{\text{EM}}$)
        have high variance across realisations.  Use $M \geq 25$ for reliable
        estimates.
\end{itemize}


% ███████████████████████████████████████████████████████████████████████████████
\chapter{Dependencies}
% ███████████████████████████████████████████████████████████████████████████████

\begin{introduction}
  \item Nine Rust crates
  \item No external SHA-256 dependency
  \item Build profile for performance
\end{introduction}

\section{Crate List}

\begin{longtable}{llp{8cm}}
  \toprule
  \textbf{Crate} & \textbf{Version} & \textbf{Purpose} \\
  \midrule
  \code{nalgebra}    & 0.33   & Dense linear algebra (eigendecomposition) \\
  \code{sprs}        & 0.11   & Sparse matrix operations ($A^2$ for Tier 1 Hasse) \\
  \code{rand}        & 0.8    & Pseudorandom number generation (\code{StdRng}) \\
  \code{rayon}       & 1.10   & Data parallelism (parallel iterators) \\
  \code{sysinfo}     & 0.33   & System memory detection \\
  \code{smallvec}    & 1.15.1 & Stack-allocated small vectors \\
  \code{rustc-hash}  & 2      & Fast non-cryptographic hashing \\
  \code{serde}       & 1      & Serialisation framework (with \code{derive}) \\
  \code{bincode}     & 1.3    & Binary serialisation for animation export \\
  \bottomrule
\end{longtable}

\section{Build Profile}

\begin{verbatim}
  [profile.release]
  opt-level = 3
  lto = "thin"
  codegen-units = 1
\end{verbatim}

This maximises performance at the cost of longer compile times.  The
\code{codegen-units = 1} setting enables whole-program optimisation.

\section{Notable Absences}

\begin{itemize}
  \item \textbf{No SHA-256 crate}: the 204-line custom implementation in
        \file{provenance.rs} avoids any external cryptographic dependency.
  \item \textbf{No TOML parser}: the configuration reader in \file{config.rs}
        uses simple line-by-line parsing (12 lines).
  \item \textbf{No logging framework}: diagnostics use \code{println!} with
        optional file-append to \file{simulation.log}.
\end{itemize}


% ═══════════════════════════════════════════════════════════════════════════════
\backmatter
% ═══════════════════════════════════════════════════════════════════════════════

\chapter*{Colophon}
\addcontentsline{toc}{chapter}{Colophon}

This document was typeset with XeLaTeX using the Modulo Synthesis palette
and tcolorbox theorem environments, matching the styling of the pedagogic
booklet series \emph{Modulo Synthesis, Volumes I \& II}.

\medskip\noindent
Source code: \url{https://github.com/ertwro/fractal_entropic_geometrodynamic}\\
DOI: \url{https://doi.org/10.5281/zenodo.18769707}\\
License: MIT

\medskip\noindent
\textcopyright{} 2026 Juan Pablo Silva Alvarado.  All rights reserved.

\end{document}
